\section{Introduction}
A woman can receive a remittance in a wallet and withdraw it all in cash. She can also receive cash and use digital payments for other transactions. A woman who runs a business using a local wallet may still need someone else to send money across a border. These are different forms of participation. Treating them as one measure of financial inclusion makes it hard to tell which part of the transaction works and which does not.

We study this distinction in Colombia. A survey of women with international migration experience separates the reported channel for the latest international remittance, the general disposition of remittance proceeds, and whether the respondent made or received a domestic digital payment or transfer in the previous 60 days. The latter two are the quantitative outcomes. They have different denominators and do not trace the same payment. Fee visibility, cash-out conditions, perceived safety, complaint outcomes and household decision arrangements provide context; they are not combined into a welfare measure.

Interviews and focus groups with Venezuelan migrant women document a contrast the survey outcomes cannot capture on their own. Women described using wallets and transfers for purchases, work and family payments within Colombia. For some transactions involving Venezuela, the local payment was instead the first step in a chain involving an intermediary, currency conversion and a separate confirmation from the recipient. Women also described different experiences in other corridors. The study therefore distinguishes digital ability, connection to a domestic payment service and the ability to complete an international transaction on terms the user can understand and verify.

This question complements work on earlier points of entry. \citet{barboni2025} study how perceived social acceptance shapes regularized Venezuelan migrants' willingness to engage with financial providers in Colombia. \citet{martinshields2022} distinguish access to technology from the use of digital government services in Bogotá. \citet{kabeer1999} separates resources, agency and achievements. We likewise distinguish access to a service from what people report doing with money and from who can decide how to use it. Neither a digital transaction nor a reported decision arrangement establishes a gain in agency.

Research on shared public digital systems distinguishes the functions of identity, payments and data exchange from their governance and adoption \citep{eavesrao2025}. Remittance studies likewise caution against reducing a transfer to a single rail: \citet{cirolia2022} describe how formal, informal and digital arrangements intersect in two Cape Town cases serving Congolese migrants. These studies do not tell us how common a specific practice is in Colombia. They help frame three questions for our evidence: what women say they do after receipt; which conditions accompany domestic use and retaining value digitally; and what can the migrant and institutional accounts tell us about where a transaction becomes difficult?

We first set the survey and fieldwork in Colombia's payments and documentation context, then separate the two engagement outcomes and their analytical bases. We read the migrant accounts before the expert diagnoses and regional cases. The discussion asks what a policy evaluation would need to measure beyond use.
